% Pista geo-25j-q4b · Geometria
% Procedència: PAU Catalunya 2025 · Convocatòria de juny · Sèrie 1 · Exercici 4 — Opció B
\documentclass[11pt,a4paper]{article}
\usepackage{pista}
\pistainfo{Catalunya 2025}{Convocatòria de juny}{Sèrie 1}

\begin{document}

\section*{\textcolor{blauPista}{Exercici 4 --- Opció B}}

\noindent Considereu el pla $\pi$ d'equació $x + y = 0$.

\subsection*{\textit{a)} \normalfont Calculeu l'equació del pla $\pi'$ que és perpendicular a $\pi$ i conté els punts $P = (1, -1, 2)$ i $Q = (3, -3, 6)$. \hspace*{\fill}\textnormal{[1~punt]}}

\pista{Com que el pla $\pi'$ és perpendicular a $\pi$, el vector normal de $\pi$, $\vec{n} = (1, 1, 0)$, és un vector director de $\pi'$. Fixa't que a l'equació de $\pi$ no hi ha $z$, i per això la 3a coordenada de $\vec{n}$ és $0$. L'altre vector director de $\pi'$ és $\vec{PQ}$.

\medskip
A continuació, escriu l'equació del pla que passa pel punt $P$ amb aquests dos vectors directors, igualant a zero aquest determinant:
\[
\begin{vmatrix} x - 1 & y + 1 & z - 2 \\ \ldots & \ldots & \ldots \\ \ldots & \ldots & \ldots \end{vmatrix} = 0
\]
Finalment, desenvolupa'l i simplifica l'equació que obtinguis.}

\subsection*{\textit{b)} \normalfont Calculeu l'equació paramètrica de la recta continguda en $\pi'$ i que conté els punts de $\pi'$ a la mateixa distància de $P$ que de $Q$. \hspace*{\fill}\textnormal{[1,5~punts]}}

\pista{Els punts que són a la mateixa distància de $P$ i de $Q$ sempre formen un pla. Per trobar aquest pla, cal saber que és un pla perpendicular al segment $PQ$ que passa pel seu punt mitjà. Per tant, el seu vector normal és $\vec{PQ}$.

\medskip
A continuació, la recta que busques és la intersecció d'aquest pla amb el pla $\pi'$ de l'apartat \textit{a)}. Resol el sistema format per les dues equacions: escriu simplement \mbox{$y = t$} (és un paràmetre lliure), i aïlla $x$ i $z$ en funció de $t$.}

\end{document}
